% =====================================================================
%  colors.tex — الألوان ومعانيها
%  كل قارئ له لون ثابت في الملف كله، وراوياه يأخذان لونه نفسه.
%  لتغيير لون قارئ غيّر القيمة هنا فقط فيتغيّر في كل المستند.
% =====================================================================

% ---------- ألوان القراء السبعة ----------
\definecolor{cNafi}{HTML}{1565C0}     % نافع وراوياه (قالون، ورش)      — أزرق
\definecolor{cIbnKathir}{HTML}{2E7D32}% ابن كثير وراوياه (البزي، قنبل) — أخضر
\definecolor{cAbuAmr}{HTML}{D35400}   % أبو عمرو وراوياه (الدوري، السوسي) — برتقالي
\definecolor{cIbnAmir}{HTML}{6A1B9A}  % ابن عامر وراوياه (هشام، ابن ذكوان) — بنفسجي
\definecolor{cAsim}{HTML}{C62828}     % عاصم وراوياه (شعبة، حفص)        — أحمر
\definecolor{cHamza}{HTML}{00838F}    % حمزة وراوياه (خلف، خلاد)        — فيروزي
\definecolor{cKisai}{HTML}{AD1457}    % الكسائي وراوياه (أبو الحارث، الدوري) — وردي غامق

% ---------- رموز الجماعة (ث خ ذ ظ غ ش، صحبة، عم ...) ----------
\definecolor{cGroup}{HTML}{5D4037}    % بني

% ---------- عناصر الشرح ----------
\definecolor{gold}{HTML}{B8860B}       % الإطارات والزخارف
\definecolor{parchment}{HTML}{FBF6E9}  % خلفية صندوق البيت
\definecolor{maroon}{HTML}{7B1F1F}     % العناوين
\definecolor{headcolor}{HTML}{4E342E}  % رأس الصفحة والعناوين الفرعية
\definecolor{wordbg}{HTML}{FFF0A8}     % تظليل الكلمة القرآنية المختلَف فيها داخل البيت
\definecolor{ayahbracket}{HTML}{B8860B}% أقواس الآيات ﴿ ﴾
\definecolor{vocabframe}{HTML}{546E7A} % صندوق معاني الكلمات
\definecolor{vocabback}{HTML}{F1F5F7}
\definecolor{ruleframe}{HTML}{2E7D32}  % صندوق القاعدة
\definecolor{ruleback}{HTML}{EEF7EE}
\definecolor{exframe}{HTML}{1565C0}    % صندوق الأمثلة
\definecolor{exback}{HTML}{EEF4FB}
\definecolor{tahrirframe}{HTML}{6A1B9A}% صندوق التحريرات (إتحاف البرية)
\definecolor{tahrirback}{HTML}{F5EEF8}
\definecolor{noteframe}{HTML}{E65100}  % صندوق التنبيه
\definecolor{noteback}{HTML}{FFF4E5}
\definecolor{sumframe}{HTML}{7B1F1F}   % صندوق الخلاصة
\definecolor{sumback}{HTML}{FDF2F2}
